\section{Översikt: vad en investmentbank gör och var värdering kommer in}
\label{sec:oversikt}

Du kan redan räkna. Det du behöver lära dig är vad som ska räknas, varför, och hur man pratar om
det. Den här guiden går igenom de fyra värderingsmetoder som en junior på en investmentbank
förväntas kunna, och den gör det med ett riktigt program som facit: \kod{valuation-engine}, som
värderar alla svenska bolag på Yahoo Finance. Varje påstående om vad programmet gör är
kontrollerat mot koden, och alla tal om Atlas Copco kommer från en och samma körning.

\subsection{Vad en investmentbank gör}

En investmentbank \eng{investment bank} tjänar pengar på rådgivning och kapitalanskaffning, inte
på in- och utlåning. Verksamheten brukar delas upp så här:

\begin{description}
  \item[Företagsaffärer] \eng{mergers and acquisitions, M\&A}. Banken ger råd till den som säljer
    ett bolag (säljuppdrag, \en{sell-side}) eller den som köper (köpuppdrag, \en{buy-side}).
    Arvodet är oftast en andel av affärens värde.
  \item[Aktiemarknad] \eng{equity capital markets, ECM}. Börsintroduktioner \eng{initial public
    offering, IPO} och nyemissioner.
  \item[Skuldmarknad] \eng{debt capital markets, DCM}. Obligationer och lån, bland annat
    finansieringen av lånefinansierade uppköp \eng{leveraged finance}.
\end{description}

Som analytiker \eng{analyst}, den första rollen efter examen, bygger du modeller i Excel och
sidor i PowerPoint. Nästan allt du gör vilar på en fråga: \emph{vad är bolaget värt, och varför?}

\subsection{Var värderingen används}

\begin{description}
  \item[Pitch.] Banken försöker vinna ett uppdrag och visar i en presentation \eng{pitch book} vad
    den tror att bolaget är värt. Slutbilden är ofta ett diagram med ett värdeintervall per
    metod, en så kallad \en{football field} (figur~\ref{fig:football} på
    sidan~\pageref{fig:football}).
  \item[Fairness opinion.] Ett skriftligt utlåtande till en styrelse om att ett bud är skäligt ur
    finansiell synvinkel. Här ska varje antagande tåla granskning.
  \item[M\&A.] Köparen vill veta vad den högst kan betala, säljaren vad den minst bör godta.
    Diskonterade kassaflöden ger ett fristående värde, jämförelsebolag visar vad marknaden
    betalar i dag.
  \item[LBO.] När köparen är ett riskkapitalbolag \eng{private equity} räknar man baklänges: vilket
    pris ger ägaren den avkastning den kräver? Det priset fungerar ofta som ett golv i
    budgivningen, eftersom en finansiell köpare saknar synergier.
\end{description}

Ingen metod ger ett sant värde. Poängen med att använda flera är att se \emph{var} de skiljer sig
och kunna förklara varför. Det är också vad kapitel~\ref{sec:case} handlar om: metoderna säger
116 till 178\kr{} per aktie för Atlas Copco när kursen är 192,25\kr, och det intressanta är
förklaringen.

\subsection{Projektet mot verkligheten}

Tabell~\ref{tab:mappning} visar vad varje del av projektet motsvarar på en bank, och var
projektet förenklar. Läs högerkolumnen noga: det är de sakerna en intervjuare frågar om.

\begin{table}[ht]
\centering\small
\caption{Projektets delar och vad de motsvarar i arbetet på en bank.}
\label{tab:mappning}
\begin{tabularx}{\linewidth}{@{}l>{\raggedright\arraybackslash}X>{\raggedright\arraybackslash}X@{}}
\toprule
Del av projektet & Motsvarar på banken & Viktigaste förenklingen \\
\midrule
\kod{data/} & Att hämta bokslut och kurser ur årsredovisningar och betalda databaser och
  justera dem för hand & Gratis data från Yahoo Finance, mekanisk mappning, senaste räkenskapsår
  i stället för rullande tolv månader \\
\kod{models/dcf.py} & DCF-modellen i Excel & Prognosen räknas mekaniskt ur tre till fyra års
  historik, ingen analytiker har läst bolagets rapporter \\
\kod{models/comps.py} & Jämförelsebolagsanalys \eng{trading comps} & Jämförelsebolagen väljs
  av en algoritm, inte av någon som kan branschen \\
\kod{models/lbo.py} & LBO-modellen & Årsmodell med två lånetrancher och schablonvillkor \\
\kod{models/monte_carlo.py} & Känslighetsanalysen. Banker visar oftast tabeller, simulering är
  ovanligare i pitchmaterial & Fyra drivare, antagna standardavvikelser \\
\kod{viz/} & Diagrammen i presentationen & Samma diagram för alla bolag \\
\kod{export/excel.py} & Själva modellen, som på en bank lever i Excel & En export med
  färdiga tal, bara ett fåtal levande formler \\
\kod{batch.py} & Screening av många bolag & Ingen manuell kontroll per bolag \\
\bottomrule
\end{tabularx}
\end{table}

\begin{komihag}
Motorn är konsekvent, inte klok. Den gör samma sak med alla 687 bolag och kraschar inte, men
den vet ingenting om bolagen. Det du ska lära dig av den är mekaniken och var omdömet behövs.
\end{komihag}

\subsection{Beteckningar}

Samma symboler används i hela guiden. Belopp anges i miljarder kronor (mdkr) om inget annat
sägs, och räntor som decimaltal i formler och procent i text.

\begin{table}[ht]
\centering\small
\caption{Beteckningar.}
\label{tab:beteckningar}
\begin{tabularx}{\linewidth}{@{}lX@{\hspace{1.5em}}lX@{}}
\toprule
Symbol & Betydelse & Symbol & Betydelse \\
\midrule
$t$ & år, $t = 1, \dots, N$ & $r$ & diskonteringsränta, i DCF lika med $\WACC$ \\
$N$ & antal prognosår (5 i motorn) & $g$ & evig tillväxt \eng{terminal growth} \\
$R_t$ & omsättning år $t$ & $r_f$ & riskfri ränta \\
$m$ & EBITDA-marginal & $\ERP$ & aktiemarknadens riskpremie \\
$d$, $c$ & avskrivningar respektive capex som andel av omsättningen & $\beta_L$, $\beta_U$ &
  beta med respektive utan skuldsättning \\
$\nu$ & rörelsekapital som andel av omsättningen & $r_e$, $r_d$ & avkastningskrav på eget
  kapital respektive låneränta före skatt \\
$\tau$ & skattesats (20,6\pct{} i projektet) & $E$, $D$ & marknadsvärde på eget kapital
  respektive skuld \\
$\FCFF_t$ & fritt kassaflöde till alla finansiärer & $\TV_N$ & slutvärde vid utgången av år
  $N$ \\
$\EV$ & rörelsevärde \eng{enterprise value} & $M$ & multipel, till exempel $\EV/\EBITDA$ \\
\bottomrule
\end{tabularx}
\end{table}
